% GENERATED FILE. Edit canonical lessons, then run npm run generate:books.
% canonical-source-hash: fnv1a64:684cc01620b00416
% canonical-lessons: JA-W11-small-tsu, JA-W11-te, JA-W11-ta, JA-W11-da, JA-C11-te-hand, JA-C11-mimi, JA-C11-kuchi, JA-C11-ashi, JA-C11-hana, JA-C11-kao, JA-C11-me-eye, JA-C11-body-map

\chapter{Seven Body Words, No New Guessing}
\label{ch:ja-body-seven}
\hlchaptermodality{11}

\begin{chapteropening}
\textbf{By the end of this chapter:} \emph{I can read, write, hear, and say seven basic body words, and I can decode the repair signs that were previously only sound-first exposure.}

Seven concrete body labels built only from earned hiragana after four tiny lessons close every outstanding repair-script exposure.
\end{chapteropening}

\section[\texorpdfstring{\ja{っ} (small tsu)}{っ (small tsu)}]{\ja{っ} — one small held beat}

\label{lesson:JA-W11-small-tsu}



\begin{quote}
Point here and say \emph{koko}. Count one time with \emph{ichido}. Ask for the line a little more slowly, then isolate \emph{yukkuri} and \emph{itte kudasai}. Hold the silence before each second consonant; do not add a vowel.
\end{quote}

\subsection*{Script — the shape}
\hlblockfigure{\detokenize{figures/JA-W11-small-tsu-filmstrip.pdf}}{How \ja{っ} is written, stroke by stroke}

\begin{quote}
\textbf{\ja{っ}} — small \emph{tsu} — one beat of consonant closure
\end{quote}

It has the same family shape as full-sized \textbf{\ja{つ}}, but it is smaller and sits lower. Copy only the small form today. In \textbf{\ja{ゆっくり}}, it holds before \emph{k}; in \textbf{\ja{いって}}, it holds before \emph{t}.

\subsection*{Your turn}
\begin{enumerate}
\raggedright
  \item Trace \textbf{\ja{っ}} once and copy it once.
  \item Hide it; write the small sign after ten seconds.
  \item Read \textbf{\ja{ゆっくり}} and \textbf{\ja{いって}} with one silent beat at \textbf{\ja{っ}}.
\end{enumerate}

\begin{tcolorbox}[breakable,colback=teal!4,colframe=teal!35!black,title={Before you move on}]
Stop after one accurate small form.

Source: \href{https://www.unicode.org/charts/PDF/U3040.pdf}{Unicode Hiragana chart}.
\end{tcolorbox}
\section[\texorpdfstring{\ja{て} (te)}{て (te)}]{\ja{て} — close \emph{itte}}

\label{lesson:JA-W11-te}



\begin{quote}
Write \textbf{\ja{もうすこし}}, say \emph{itte kudasai}, then write the small \textbf{\ja{っ}} that holds before \emph{t}.
\end{quote}

\subsection*{Script — the shape}
\hlblockfigure{\detokenize{figures/JA-W11-te-filmstrip.pdf}}{How \ja{て} is written, stroke by stroke}

\begin{quote}
\textbf{\ja{て}} — \emph{te}
\end{quote}

One continuous stroke: begin across the top, turn back, then sweep down and right. Keep the lower curve open.

\subsection*{Your turn}
\begin{enumerate}
\raggedright
  \item Trace \textbf{\ja{て}} once.
  \item Copy it once, then retrieve it after ten seconds.
  \item Join the earned ending: \textbf{\ja{い} | \ja{っ} | \ja{て}}.
\end{enumerate}

\begin{tcolorbox}[breakable,colback=teal!4,colframe=teal!35!black,title={Before you move on}]
Read \textbf{\ja{いって}} once without romanization.

Source: \href{https://www.unicode.org/charts/PDF/U3040.pdf}{Unicode Hiragana chart}; stroke order per \emph{Hitsujun Shido no Tebiki} (Japanese Ministry of Education, 1958).
\end{tcolorbox}
\section[\texorpdfstring{\ja{た} (ta)}{た (ta)}]{\ja{た} — close \emph{wakarimashita}}

\label{lesson:JA-W11-ta}



\begin{quote}
Say \emph{wakarimashita}. Write \textbf{\ja{と}} and notice that today's \emph{ta} belongs to the same consonant row but has its own shape.
\end{quote}

\subsection*{Script — the shape}
\hlblockfigure{\detokenize{figures/JA-W11-ta-filmstrip.pdf}}{How \ja{た} is written, stroke by stroke}

\begin{quote}
\textbf{\ja{た}} — \emph{ta}
\end{quote}

Four strokes: a short top stroke, a crossing vertical curve, then two small right-side strokes. Keep the right pair separate.

\subsection*{Your turn}
\begin{enumerate}
\raggedright
  \item Trace \textbf{\ja{た}} once and copy it once.
  \item Hide it; write it after ten seconds.
  \item Add it to the known final block \textbf{\ja{し} | \ja{た}}.
\end{enumerate}

\begin{tcolorbox}[breakable,colback=teal!4,colframe=teal!35!black,title={Before you move on}]
Read the ending \textbf{\ja{した}} once.

Source: \href{https://www.unicode.org/charts/PDF/U3040.pdf}{Unicode Hiragana chart}; stroke order per \emph{Hitsujun Shido no Tebiki} (Japanese Ministry of Education, 1958).
\end{tcolorbox}
\section[\texorpdfstring{\ja{だ} (da)}{だ (da)}]{\ja{だ} — voice the known base}

\label{lesson:JA-W11-da}



\begin{quote}
Write \textbf{\ja{た}}, then add the two-stroke dakuten without changing the base. Say the \emph{da} in \emph{kudasai}.
\end{quote}

\subsection*{Script — the change}
\hlblockfigure{\detokenize{figures/JA-W11-da-filmstrip.pdf}}{How \ja{だ} is written, stroke by stroke}

\begin{quote}
\textbf{\ja{た}} + dakuten $\to$ \textbf{\ja{だ}} — \emph{ta} $\to$ \emph{da}
\end{quote}

The voicing mark does the same job it did before. This is transfer, not a new complex shape.

\subsection*{Your turn}
\begin{enumerate}
\raggedright
  \item Write \textbf{\ja{た}} once.
  \item Add dakuten to make \textbf{\ja{だ}}.
  \item Point to \textbf{\ja{だ}} while saying the middle of \emph{kudasai}.
\end{enumerate}

\begin{tcolorbox}[breakable,colback=teal!4,colframe=teal!35!black,title={Before you move on}]
Read \textbf{\ja{ください}} once, attending to \textbf{\ja{だ}}.

Source: \href{https://www.unicode.org/charts/PDF/U3040.pdf}{Unicode Hiragana chart}.
\end{tcolorbox}
\section[\texorpdfstring{\ja{て} (te)}{て (te)}]{\ja{て} — hand}

\label{lesson:JA-C11-te-hand}



\begin{quote}
\emph{Write it:} \textbf{\ja{て}} and \textbf{\ja{だ}}

Picture the spot where you are and say \emph{koko}, then ask for the line a little more slowly.
\end{quote}

\subsection*{What to know first}
\begin{quote}
\textbf{\ja{て}} — \emph{te} — hand
\end{quote}

The one-sign word names the hand. Sound, meaning, reading, and writing can meet without adding another shape.

\subsection*{Your turn}
\begin{enumerate}
\raggedright
  \item Hear \emph{te}; say what it names — a hand.
  \item \emph{Read it:} \textbf{\ja{て}}, then say “hand.”
  \item \emph{Write it:} its one sign, with the word hidden
\end{enumerate}

\begin{tcolorbox}[breakable,colback=teal!4,colframe=teal!35!black,title={Before you move on}]
Stop after one clean retrieval.

Source: \href{https://words.marugotoweb.jp/}{Japan Foundation Marugoto A1 vocabulary database}.
\end{tcolorbox}
\section[\texorpdfstring{\ja{みみ} (mimi)}{みみ (mimi)}]{\ja{みみ} — ear}

\label{lesson:JA-C11-mimi}



\begin{quote}
\emph{Write it:} \textbf{\ja{み}} once

Then retrieve \textbf{\ja{て}} — hand.
\end{quote}

\subsection*{What to know first}
\begin{quote}
\textbf{\ja{みみ}} — \emph{mimi} — ear; ears
\end{quote}

No new sign: repeat \textbf{\ja{み}} on two even morae.

\subsection*{Your turn}
\begin{enumerate}
\raggedright
  \item Hear \emph{mimi}; say what it names — an ear.
  \item \emph{Read it:} \textbf{\ja{み} | \ja{み}} aloud
  \item \emph{Write it:} the repeated sign twice, with the word hidden
\end{enumerate}

\begin{tcolorbox}[breakable,colback=teal!4,colframe=teal!35!black,title={Before you move on}]
Keep the two morae even.

Source: \href{https://words.marugotoweb.jp/}{Japan Foundation Marugoto A1 vocabulary database}.
\end{tcolorbox}
\section[\texorpdfstring{\ja{口} (kuchi)}{口 (kuchi)}]{\ja{口} — mouth}

\label{lesson:JA-C11-kuchi}



\begin{quote}
\emph{Write it:} \textbf{\ja{く}} and \textbf{\ja{ち}} separately

Retrieve \textbf{\ja{みみ}}.
\end{quote}

\subsection*{What to know first}
\begin{quote}
\textbf{\ja{口}} (\ja{くち}) — \emph{kuchi} — mouth
\end{quote}

The three-stroke box you have written since chapter five is the word: \textbf{\ja{口}} is \emph{kuchi}. Its reading joins two earned signs, \textbf{\ja{く}} and \textbf{\ja{ち}}; keep two even morae.

\subsection*{Your turn}
\begin{enumerate}
\raggedright
  \item Hear \emph{kuchi}; say what it names — the mouth.
  \item \emph{Read it:} \textbf{\ja{口}}, then its reading \textbf{\ja{く} | \ja{ち}}
  \item \emph{Write it:} \textbf{\ja{口}}, then the two signs in order, with the word hidden
\end{enumerate}

\begin{tcolorbox}[breakable,colback=teal!4,colframe=teal!35!black,title={Before you move on}]
Contrast \emph{mimi} and \emph{kuchi} once.

Sources: \href{https://words.marugotoweb.jp/}{Japan Foundation Marugoto A1 vocabulary database}; \href{https://kanjivg.tagaini.net/kanjivg/kanji/053e3.html}{KanjiVG stroke-order data}.
\end{tcolorbox}
\section[\texorpdfstring{\ja{あし} (ashi)}{あし (ashi)}]{\ja{あし} — foot or leg}

\label{lesson:JA-C11-ashi}



\begin{quote}
\emph{Write it:} \textbf{\ja{あ}} and \textbf{\ja{し}}

Retrieve \textbf{\ja{くち}} — mouth.
\end{quote}

\subsection*{What to know first}
\begin{quote}
\textbf{\ja{あし}} — \emph{ashi} — foot; leg
\end{quote}

Context chooses the narrower English translation. The Japanese word stays the same.

\subsection*{Your turn}
\begin{enumerate}
\raggedright
  \item Hear \emph{ashi}; say what it names — a foot or leg.
  \item \emph{Read it:} \textbf{\ja{あ} | \ja{し}}
  \item \emph{Write it:} both known signs, with the word hidden
\end{enumerate}

\begin{tcolorbox}[breakable,colback=teal!4,colframe=teal!35!black,title={Before you move on}]
Stop after one accurate retrieval.

Source: \href{https://words.marugotoweb.jp/}{Japan Foundation Marugoto A1 vocabulary database}.
\end{tcolorbox}
\section[\texorpdfstring{\ja{はな} (hana)}{はな (hana)}]{\ja{はな} — nose}

\label{lesson:JA-C11-hana}



\begin{quote}
\emph{Write it:} \textbf{\ja{は}} and \textbf{\ja{な}}

Retrieve \textbf{\ja{あし}}.
\end{quote}

\subsection*{What to know first}
\begin{quote}
\textbf{\ja{はな}} — \emph{hana} — nose
\end{quote}

Inside this word, \textbf{\ja{は}} is \emph{ha}. Its \emph{wa} reading was tied to the topic marker job, not to every appearance of the sign.

\subsection*{Your turn}
\begin{enumerate}
\raggedright
  \item Hear \emph{hana}; say what it names — the nose.
  \item \emph{Read it:} \textbf{\ja{は} | \ja{な}} as \emph{ha-na}
  \item \emph{Write it:} both signs, with the word hidden
\end{enumerate}

\begin{tcolorbox}[breakable,colback=teal!4,colframe=teal!35!black,title={Before you move on}]
Say \emph{hana} once without turning the first mora into \emph{wa}.

Source: \href{https://words.marugotoweb.jp/}{Japan Foundation Marugoto A1 vocabulary database}.
\end{tcolorbox}
\section[\texorpdfstring{\ja{かお} (kao)}{かお (kao)}]{\ja{かお} — face}

\label{lesson:JA-C11-kao}



\begin{quote}
\emph{Write it:} \textbf{\ja{か}} and \textbf{\ja{お}}

Retrieve \textbf{\ja{はな}} — nose.
\end{quote}

\subsection*{What to know first}
\begin{quote}
\textbf{\ja{かお}} — \emph{kao} — face
\end{quote}

Give \textbf{\ja{か}} and \textbf{\ja{お}} one mora each: \emph{ka-o}.

\subsection*{Your turn}
\begin{enumerate}
\raggedright
  \item Hear two beats, \emph{ka-o}; say what it names — the face.
  \item \emph{Read it:} \textbf{\ja{か} | \ja{お}}
  \item \emph{Write it:} the two signs, with the word hidden
\end{enumerate}

\begin{tcolorbox}[breakable,colback=teal!4,colframe=teal!35!black,title={Before you move on}]
Contrast \textbf{\ja{はな}} inside the \textbf{\ja{かお}} once.

Source: \href{https://words.marugotoweb.jp/}{Japan Foundation Marugoto A1 vocabulary database}.
\end{tcolorbox}
\section[\texorpdfstring{\ja{め} (me)}{め (me)}]{\ja{め} — eye}

\label{lesson:JA-C11-me-eye}



\begin{quote}
\emph{Write it:} \textbf{\ja{め}}

Retrieve \textbf{\ja{かお}} and \textbf{\ja{みみ}}.
\end{quote}

\subsection*{What to know first}
\begin{quote}
\textbf{\ja{め}} — \emph{me} — eye; eyes
\end{quote}

One learned sign is the complete everyday word.

\subsection*{Your turn}
\begin{enumerate}
\raggedright
  \item Hear \emph{me}; say what it names — an eye.
  \item \emph{Read it:} \textbf{\ja{め}}, then say the meaning
  \item \emph{Write it:} the one-sign word, with the model hidden
\end{enumerate}

\begin{tcolorbox}[breakable,colback=teal!4,colframe=teal!35!black,title={Before you move on}]
Keep \textbf{\ja{め}} and \textbf{\ja{みみ}} distinct.

Source: \href{https://words.marugotoweb.jp/}{Japan Foundation Marugoto A1 vocabulary database}.
\end{tcolorbox}
\section[\texorpdfstring{\ja{て・みみ・くち} … (te, mimi, kuchi …)}{て・みみ・くち … (te, mimi, kuchi …)}]{Seven readable body words}

\label{lesson:JA-C11-body-map}



\begin{quote}
\emph{Write it:} \textbf{\ja{っ}, \ja{て}, \ja{た}, \ja{だ}}

\emph{Read it:} \textbf{\ja{ゆっくり}}, \textbf{\ja{いって}}, \textbf{\ja{わかりました}}, and \textbf{\ja{ください}} once — the old repair exposure is now closed
\end{quote}

\subsection*{What to know first}
\begin{quote}
\textbf{\ja{て} \ja{・} \ja{みみ} \ja{・} \ja{くち} \ja{・} \ja{あし} \ja{・} \ja{はな} \ja{・} \ja{かお} \ja{・} \ja{め}}
\end{quote}

This is a reading-and-writing map, not a new-word dump. Every word arrived in its own lesson, and every sign was already earned.

\subsection*{Your turn}
\begin{enumerate}
\raggedright
  \item \textbf{Listen:} say the English for the named body part.
  \item \textbf{Speak:} name any three without reading.
  \item \textbf{Reading:} \emph{Read it:} all seven from left to right
  \item \emph{Write it:} any three of the seven, from memory
\end{enumerate}

\begin{tcolorbox}[breakable,colback=teal!4,colframe=teal!35!black,title={Before you move on}]
Stop after one accurate three-word dictation.

Source: \href{https://words.marugotoweb.jp/}{Japan Foundation Marugoto A1 vocabulary database}.
\end{tcolorbox}
